\section{Sites}

We began the course by discussing ``sheaves on the category of sets'', which is certainly not a special case of $\Sh(X)$ for $X$ a topological space, since $\Set$ is not equivalent to $\Open(X)$ for any topological space $X$. In this section, we will discuss a generalization of the notion of sheaf which encompasses both sheaves on topological spaces and sheaves on the category of sets. For this, we must replace $\Open(X)$ and $\Set$ with a more general datum called a \emph{site}. A site is a category together with a collection of ``covering families'' -- these covering families specify exactly which equalizers should appear in the sheaf condition. It will turn out that every category can be equipped with a canonical choice of covering families, and that sheaves on $\Set$ or $\Open(X)$ with respect to these canonical covering families are exactly the sheaves we have already discussed.

\subsection{Grothendieck Topologies and the Sheaf Condition}

\begin{definition}
    Let $\catC$ be a category. A \emph{sieve} on an object $x \in \catC$ is a subfunctor $S \subseteq y(x)$ of the representable presheaf $y(x) = \Hom_\catC({-},x)$. That is, for all objects $c \in \catC$, we have a subset $S(c) \subseteq \Hom_\catC(c,x)$, and for all morphisms $f : c' \to c$ in $\catC$, we have $g \circ f \in S(c')$ for all $g \in S(c)$.
\end{definition}

\begin{warning}
    If $\catC$ is a large category, then $\PSh(\catC)$ is not necessarily a category in the usual sense (its hom-sets may in fact be proper classes). However, the notions of presheaf and subfunctor still make sense, so in particular it is fine to talk about sieves. I will try to be careful about size issues in what follows, but we'll be forced to return to this technical point again soon, and it is always worth thinking about which things are ``small'' and ``large''! This is a spot where the theory of Grothendieck universes would be very useful to make things precise, but I don't want to assume knowledge of Grothendieck universes in this course. I've included an appendix \Cref{appendix:size issues} which introduces Grothendieck universes, but you can safely ignore it.
\end{warning}

A sieve is, in other words, a collection of morphisms into $x$ which is closed under precomposition. It's clear that the collection of sieves on $x$ is closed under arbitrary intersection (the empty intersection is the maximum sieve $y(x)$). Thus, we may talk about the sieve on $x$ generated by a collection of morphisms $\{f_\alpha\}_\alpha$ with codomain $x$, which is the smallest sieve containing all of the $f_\alpha$. Explicitly, the sieve generated by $\{f_\alpha\}_\alpha$ consists of all morphisms which factor through some $f_\alpha$, i.e.
\[S(c) = \{g \in \catC(c,x) : g = f_\alpha \circ h \text{ for some } \alpha \text{ and some }h : c \to \dom(f_\alpha)\}.\]

\begin{example}\label{example:chap:sheaves:sec:sites:sieve generated by open cover}
    Let $X$ be a topological space, and let $U \subseteq X$ be an open set. Let $\{U_\alpha\}_\alpha$ be an open cover of $U$. The sieve on $U \in \Open(X)$ generated by the inclusions $U_\alpha \hookrightarrow U$ is the subfunctor $S \subseteq y(U)$ defined by
    \[S(V) = \{V \hookrightarrow U : V \subseteq U_\alpha \text{ for some } \alpha\}.\]
\end{example}

If $S$ is a sieve on $x$ and $f : x' \to x$ is a morphism in $\catC$, we may define the \emph{pullback sieve} $f^* S$ on $x'$ by
\[(f^* S)(z) = \{g \in \catC(z,x') : f \circ g \in S(z)\}.\]

\begin{example}
    Let $X$ be a topological space, let $U \subseteq X$ be an open set, and let $\{U_\alpha\}_\alpha$ be an open cover of $U$. Let $S \subseteq y(U)$ be the sieve generated by the inclusions $U_\alpha \hookrightarrow U$. Let $V \subseteq U$ be an open subset. Then the pullback sieve $V^* S$ on $V$ is the sieve generated by the inclusions $V \cap U_\alpha \hookrightarrow V$, i.e.
    \[(V^* S)(W) = \{W \hookrightarrow V : W \subseteq V \cap U_\alpha \text{ for some } \alpha\}.\]
\end{example}

\begin{definition}
    Let $\catC$ be a category. A \emph{Grothendieck topology} $\tau$ on $\catC$ is a choice of a collection\footnote{Potentially a proper class!} of sieves $\tau(x)$ on each object $x \in \catC$, called the \emph{covering sieves}, such that:
    \begin{enumerate}[label=(T\arabic*)]
        \item For all $S \in \tau(x)$ and all morphisms $f : y \to x$, $f^* S \in \tau(y)$.
        \item For all $S \in \tau(x)$ and all sieves $T$ on $x$, if $f^* T \in \tau(y)$ for all $f : y \to x$ in $S$, then $T \in \tau(x)$.
        \item $y(x) \in \tau(x)$ for all $x \in \catC$.
    \end{enumerate}
    A \emph{site} is a pair $(\catC, \tau)$ of a category $\catC$ and a Grothendieck topology $\tau$ on $\catC$.
\end{definition}

\startlecture[Sep 24, 2026]

\begin{example}\label{example:canonical topology on OpenX}
    Let $X$ be a topological space. There is a Grothendieck topology $\tau$ on $\Open(X)$ defined by declaring the covering sieves on $U$ to be precisely the sieves generated by open covers of $U$, as in \Cref{example:chap:sheaves:sec:sites:sieve generated by open cover}.
\end{example}

\begin{example}\label{example:canonical topology on Set}
    There is a Grothendieck topology on $\Set$ defined by declaring the covering sieves $S \in \tau(X)$ to be precisely those for which the morphisms in $S$ are jointly surjective, i.e.\ for all $x \in X$ there exists some $f : Y \to X$ in $S(Y)$ and some $y \in Y$ such that $f(y) = x$.
\end{example}

\begin{remark}
    One sometimes sees sites defined in a slightly different way, in terms of \emph{Grothendieck pre-topologies}. Here, instead of specifying covering sieves, one specifies covering families $\{f_\alpha : x_\alpha \to x\}_\alpha$ of morphisms with codomain $x$. These covering families are required to satisfy the following axioms:
    \begin{enumerate}[label=(PT\arabic*)]
        \item For all covering families $\{f_\alpha : x_\alpha \to x\}_\alpha$ and all morphisms $g : y \to x$, $\{x_\alpha \times_x y \to y\}_\alpha$ is a covering family.
        \item For all covering families $\{f_\alpha : x_\alpha \to x\}_\alpha$ and all collections of covering families $\{g_{\alpha,\beta} : x_{\alpha,\beta} \to x_\alpha\}_\beta$, the family $\{f_\alpha \circ g_{\alpha,\beta} : x_{\alpha,\beta} \to x\}_{\alpha,\beta}$ is a covering family.
        \item The singleton family $\{\id_x : x \to x\}$ is a covering family for all $x$.
    \end{enumerate}
    Any Grothendieck pre-topology generates a Grothendieck topology by declaring a sieve $S$ on $x$ to be covering if and only if it contains a covering family\footnote{It is worth taking the time to understand how the axioms PT1-PT3 translate to the axioms T1-T3 via this construction!}. There are two downsides to this approach:
    \begin{itemize}
        \item Many distinct pre-topologies may generate the same topology, and the eventual definition of ``sheaves on a site'' that we will use only depends on the Grothendieck topology, not the pre-topology.
        \item The axioms of a Grothendieck pre-topology assume the existence of (at least some) pullbacks in $\catC$, while the axioms of a Grothendieck topology do not.
    \end{itemize}
\end{remark}

Suppose for a moment that $\catC$ is small. The Yoneda lemma says that
\[\cal{F}(x) \cong \Hom_{\PSh(\catC)}(y(x),\cal{F})\]
for all presheaves $\cal{F} \in \PSh(\catC)$. Consequently, for any sieve $S \subseteq y(x)$, we have a canonical map
\[\cal{F}(x) \cong \Hom_{\PSh(\catC)}(y(x),\cal{F}) \to \Hom_{\PSh(\catC)}(S,\cal{F}).\]
We say that $\cal{F}$ \emph{satisfies the sheaf condition with respect to $S$} or \emph{is local with respect to $S$} if this map is a bijection. We will write $\cal{F}(S)$ as shorthand for $\Hom_{\PSh(\catC)}(S,\cal{F})$.

If $\catC$ is not a small category, the version of the Yoneda lemma we have stated (\Cref{prop:Yoneda}) does not apply. However, much of it still holds:
\begin{itemize}
    \item It is still true that $\Hom(y(x),\cal{F}) \cong \cal{F}(x)$ naturally in presheaves $\cal{F}$ and objects $x$; in particular, $\Hom(y(x),\cal{F})$ is a set.
    \item It is still true that every presheaf $\cal{F}$ is a colimit of representable presheaves in a canonical way (this colimit is indexed by a potentially large category, but it does exist).
\end{itemize}
In this case, we can still obtain a canonical map $\cal{F}(x) \to \cal{F}(S) := \Hom(S,\cal{F})$ for any sieve $S$ on $x$, but $\cal{F}(S)$ may be a proper class. Still, it makes sense to ask whether or not $\cal{F}$ satisfies the sheaf condition with respect to $S$!

If $S$ is generated by a set of morphisms, then we can guarantee that $\cal{F}(S)$ is a set. Many sites encountered in practice have the property that all covering sieves are generated by sets.

\begin{proposition}\label{prop:chap:sheaves:sec:sites:FS is a set for small sieves}
    Let $S$ be a sieve on $x \in \catC$. If $S$ is generated by a set of morphisms $\{f_i : c_i \to x\}_i$, then
    \[\cal{F}(S) := \Hom(S,\cal{F})\]
    is a set for any presheaf $\cal{F} \in \PSh(\catC)$.
\end{proposition}
\begin{proof}
    The key is that, in this case, the morphism
    \[\coprod_i y(c_i) \to S\]
    determined by the elements $f_i \in S(c_i)$ is an epimorphism in $\PSh(\catC)$; i.e.\ it is a natural transformation whose components are all surjections. To see this, note that the image of this map $\coprod_i y(c_i) \to S$ is a well-defined subfunctor of $S$ (thus in particular a sieve on $x$) which contains all of the $f_i$; since $S$ is the smallest sieve containing all of the $f_i$, it follows that this image must be all of $S$. It follows that the induced map
    \[\cal{F}(S) = \Hom(S,\cal{F}) \to \Hom(\coprod_i y(c_i),\cal{F}) \cong \prod_i \Hom(y(c_i),\cal{F}) \cong \prod_i \cal{F}(c_i)\]
    is injective, and thus $\Hom(S,\cal{F})$ is a set.
\end{proof}

\begin{proposition}\label{prop:chap:sheaves:sec:sites:sheaf condition as limit}
    Let $\catC$ be a category, and let $S$ be a sieve on $x \in \catC$. Let $S/x$ denote the full subcategory of $\catC/x$ spanned by the morphisms in $S$.

    Then, for any presheaf $\cal{F}$, elements of $\cal{F}(S)$ are in bijection with compatible families $(s_f)_{f \in S}$, where $s_f \in \cal{F}(c)$ for each $f : c \to x$ in $S$. Here compatibility means that for all morphisms $g : f' \to f$ in $S/x$, we have $\cal{F}(g)(s_f) = s_{f'}$.
\end{proposition}
\begin{proof}
    Given $\alpha : S \to \cal{F}$, set $s_f := \alpha_c(f) \in \cal{F}(c)$ for each $f : c \to x$ in $S$. This is compatible by naturality of $\alpha$. In the other direction, given a compatible family $(s_f)_{f \in S}$, define $\alpha_c : S(c) \to \cal{F}(c)$ by $\alpha_c(f) := s_f$. These are the components of a natural transformation $\alpha : S \to \cal{F}$ by the compatibility condition.
\end{proof}

This is almost just saying that $\cal{F}(S)$ is the limit of the diagram
\[(\smallint S)^\op \cong (S/x)^\op \xrightarrow{\dom} \catC^\op \xrightarrow{\cal{F}} \Set,\]
coming from the fact that $S$ is the colimit of the diagram
\[\smallint S \cong S/x \xrightarrow{\dom} \catC \xrightarrow{y} \PSh(\catC).\]
The only trouble is that the indexing category $\int S$ is potentially large, and so the desired limit does not necessarily exist in $\Set$, as we have discussed. However, the colimit defining $S$ does exist in $\PSh$, and we can describe directly what elements of $\cal{F}(S) = \Hom(S,\cal{F})$ are by the universal property of colimits. This gives the statement in the proposition.

\begin{proposition}\label{prop:chap:sheaves:sec:sites:sheaf condition as equalizer}
    Let $\catC$ be a category, and let $S$ be a sieve on $x \in \catC$. Assume that the slice category $\catC/x$ admits binary products\footnote{These are just pullback squares with $x$ in the corner.} and that $S$ is generated by a set of morphisms $\{f_i : c_i \to x\}_i$. Then $\cal{F}(S)$ is naturally isomorphic to the equalizer of the diagram
    \[\prod_{i} \cal{F}(c_i) \rightrightarrows \prod_{i,j} \cal{F}(c_i \times_x c_j).\]
\end{proposition}
\begin{proof}
    Recall from the proof of \Cref{prop:chap:sheaves:sec:sites:FS is a set for small sieves} that restriction along the $f_i$ gives an injection $\cal{F}(S) \hookrightarrow \prod_i \cal{F}(c_i)$; we must identify its image with the equalizer, and for this we use the description of \Cref{prop:chap:sheaves:sec:sites:sheaf condition as limit}.

    If $(s_f)_{f \in S}$ is a compatible family, then for all $i,j$ we have $f_i \circ p_1 = f_j \circ p_2$, where $p_1 : c_i \times_x c_j \to c_i$ and $p_2 : c_i \times_x c_j \to c_j$ are the canonical projections, so
    \[\cal{F}(p_1)(s_{f_i}) = s_{f_i \circ p_1} = s_{f_j \circ p_2} = \cal{F}(p_2)(s_{f_j}).\]

    Conversely, suppose $(s_i)_i$ satisfies $\cal{F}(p_1)(s_i) = \cal{F}(p_2)(s_j)$ for all $i,j$. Every $f : d \to x$ in $S$ factors as $f = f_i \circ h$ for some $i$ and some $h : d \to c_i$; set $s_f := \cal{F}(h)(s_i)$. This is independent of the factorization: if also $f = f_j \circ k$, then $(h,k)$ induces $r : d \to c_i \times_x c_j$ with $h = p_1 \circ r$ and $k = p_2 \circ r$, whence
    \[\cal{F}(h)(s_i) = \cal{F}(r)(\cal{F}(p_1)(s_i)) = \cal{F}(r)(\cal{F}(p_2)(s_j)) = \cal{F}(k)(s_j).\]
    The family $(s_f)_{f \in S}$ is compatible: if $g : f' \to f$ is a morphism in $S/x$ and $f = f_i \circ h$, then $f' = f_i \circ (h \circ g)$, so $s_{f'} = \cal{F}(h \circ g)(s_i) = \cal{F}(g)(s_f)$. Finally, taking $h = \id_{c_i}$ shows $s_{f_i} = s_i$.
\end{proof}

This understanding allows us to extend our sheaf condition so it applies to presheaves with values in any algebraic category.

\begin{definition}
    Let $\catC$ be a category, let $S$ be a sieve on $x \in \catC$, and let $\catD$ be an algebraic category. A presheaf $\cal{F} \in \PSh(\catC,\catD)$ is said to \emph{satisfy the sheaf condition with respect to $S$} (or be \emph{local with respect to $S$}) if its underlying $\Set$-valued presheaf (coming from the forgetful functor $\catD \to \Set$) is local with respect to $S$.
\end{definition}

We have same equalizer description of the sheaf condition in this more general setting, coming from the fact that limits in $\catD$ are computed in $\Set$.

\begin{proposition}\label{prop:chap:sheaves:sec:sites:sheaf condition as equalizer in algebraic category}
    Let $\catD$ be an algebraic category. If $S$ is generated by a set of morphisms $\{f_i : c_i \to x\}_i$ and $\catC/x$ admits binary products, then $\cal{F} \in \PSh(\catC,\catD)$ is local with respect to $S$ if and only if
    \[\cal{F}(x) \to \prod_i \cal{F}(c_i) \rightrightarrows \prod_{i,j} \cal{F}(c_i \times_x c_j)\]
    is an equalizer diagram in $\catD$.
\end{proposition}

This allows us to define sheaves on a site with values in an algebraic category\footnote{Or, as before, any locally finitely presentable category, if you're willing to be more abstract.}.

\begin{definition}\label{def:chap:sheaves:sec:sites:sheaf}
    Let $(\catC, \tau)$ be a site, and let $\catD$ be an algebraic category. A presheaf $\cal{F} \in \PSh(\catC,\catD)$ is a \emph{sheaf} if it is local with respect to all covering sieves. We write $\Sh((\catC,\tau),\catD)$, or alternatively $\Sh_{\tau}(\catC,\catD)$, for the full subcategory of $\PSh(\catC,\catD)$ spanned by the sheaves. When $\catD = \Set$, we omit it from the notation, writing $\Sh(\catC,\tau)$ or $\Sh_{\tau}(\catC)$. When the topology $\tau$ is clear from context, we omit it as well, writing simply $\Sh(\catC,\catD)$ or $\Sh(\catC)$.
\end{definition}

Again, we warn that $\Sh(\catC,\catD)$ may fail to be a category in the traditional sense if $\catC$ is not small (because its hom-sets may be too large).

\subsection{The Canonical Topology}

So far, we have only considered sheaves on the category $\Set$, and on the categories $\Open(X)$ for $X$ a topological space. In both of these cases, we had a ``canonical choice'' of Grothendieck topology, which encoded our intuitions about what ``covers'' are. In fact, every category $\catC$ has a canonical choice of Grothendieck topology, called the \emph{canonical topology}. In order to understand this canonical topology, we make a few observations:


\begin{enumerate}
    \item $\PSh(\catC)$ is the free cocompletion of $\catC$. Informally, $\Sh_{\tau}(\catC)$ is a (not necessarily free) cocompletion of $\catC$: it is generated under colimits by the representable presheaves, subject to the relations that covering sieves should ``become colimits''. Nothing requires the covering sieves to be colimits in $\catC$ itself.
    \item A representable presheaf $y(c)$ is a sheaf if and only if the object $c$ ``believes'' that every covering sieve is a colimit in $\catC$. If every representable presheaf is a sheaf, then every object of $\catC$ believes that every covering sieve is a colimit in $\catC$, and this implies that every covering sieve is indeed a colimit in $\catC$.
\end{enumerate}
In this subsection, we make the second observation precise, and use it to construct the canonical topology.

\begin{definition}\label{def:chap:sheaves:sec:sites:colimit sieve}
    Let $S$ be a sieve on $x \in \catC$. We say that $S$ is a \emph{colimit sieve} if the cocone $(f : \dom f \to x)_{f \in S}$ exhibits $x$ as the colimit of the diagram $S/x \xrightarrow{\dom} \catC$. We say that $S$ is a \emph{universal colimit sieve} if the pullback sieve $g^* S$ is a colimit sieve for every morphism $g : x' \to x$.\footnote{This terminology follows Lester \cite{lester_covers_2019,lester_canonical_2019}, where these are called (universal) \emph{colim sieves}. The notion itself goes back to Verdier in SGA~4 \cite[Exp.~II]{artin_theorie_1972}; Johnstone \cite[\nopp C2.1]{johnstone_sketches_2002} calls them (universally) \emph{effective-epimorphic} sieves.}
\end{definition}

Of course, every universal colimit sieve is a colimit sieve, since we may choose $g = \id_x$.

\begin{proposition}\label{prop:chap:sheaves:sec:sites:when representables are sheaves}
    Let $S$ be a sieve on $x \in \catC$, and let $c \in \catC$. The representable presheaf $y(c)$ is local with respect to $S$ if and only if every cocone from the diagram $S/x \xrightarrow{\dom} \catC$ to $c$ factors uniquely through the cocone $(f : \dom f \to x)_{f \in S}$. Consequently, $S$ is a colimit sieve if and only if every representable presheaf is local with respect to $S$.
\end{proposition}
\begin{proof}
    By \Cref{prop:chap:sheaves:sec:sites:sheaf condition as limit}, an element of $y(c)(S) = \Hom(S,y(c))$ is a family of morphisms $(s_f : \dom f \to c)_{f \in S}$ such that $s_f \circ g = s_{f'}$ for all morphisms $g : f' \to f$ in $S/x$; this is precisely a cocone from the diagram $S/x \to \catC$ to $c$. The map $y(c)(x) \to y(c)(S)$ sends $h : x \to c$ to the cocone $(h \circ f)_{f \in S}$, which is a bijection if and only if every cocone to $c$ factors uniquely through the cocone $(f : \dom f \to x)_{f \in S}$.
\end{proof}

This makes precise the sense in which $c$ ``believes'' that $S$ is a colimit: $x$ has the universal property of the colimit of $S/x \to \catC$, as long as we only test this universal property against morphisms into $c$.

\begin{definition}\label{def:chap:sheaves:sec:sites:subcanonical}
    A Grothendieck topology $\tau$ on $\catC$ is \emph{subcanonical} if every representable presheaf is a sheaf.
\end{definition}

\begin{corollary}\label{cor:chap:sheaves:sec:sites:subcanonical iff colimit sieves}
    For a Grothendieck topology $\tau$ on $\catC$, the following are equivalent:
    \begin{enumerate}[label=(\roman*)]
        \item $\tau$ is subcanonical;
        \item every covering sieve is a colimit sieve;
        \item every covering sieve is a universal colimit sieve.
    \end{enumerate}
\end{corollary}
\begin{proof}
    Conditions (i) and (ii) are equivalent by \Cref{prop:chap:sheaves:sec:sites:when representables are sheaves}, and (iii) implies (ii) since every universal colimit sieve is a colimit sieve. Finally, (ii) implies (iii): if $S \in \tau(x)$, then $g^* S \in \tau(x')$ for every $g : x' \to x$ by (T1), so $g^* S$ is a colimit sieve by (ii).
\end{proof}

\begin{example}\label{example:chap:sheaves:sec:sites:dense topology}
    Let $X$ be a space. Declare a sieve $S$ on $U \in \Open(X)$ to be covering if the union of the open sets in $S$ is dense in $U$. This defines a Grothendieck topology on $\Open(X)$ (\Cref{exercise:chap:sheaves:sec:sites:dense topology is a topology}), called the \emph{dense topology}, and it is not subcanonical in general. For instance, take $X = \RR$, and let $S$ be the sieve on $\RR$ generated by the inclusion $\RR \setminus \{0\} \hookrightarrow \RR$, which is covering since $\RR \setminus \{0\}$ is dense in $\RR$. The inclusions $W \hookrightarrow \RR \setminus \{0\}$ for $W \in S$ form a cocone from $S/\RR$ to $\RR \setminus \{0\}$, and this cocone cannot factor through $\RR$, since there is no morphism $\RR \to \RR \setminus \{0\}$ in $\Open(\RR)$. So $S$ is not a colimit sieve, and $y(\RR \setminus \{0\})$ is not a sheaf: the object $\RR \setminus \{0\}$ does not believe that $S$ is a colimit.
\end{example}

\begin{example}
    The empty sieve on $x$ is a colimit sieve if and only if $x$ is an initial object. So, any subcanonical topology must have the property that no empty sieve covers a non-initial object. \begin{draft}Conversely, if $x$ is an initial object, then the empty sieve covers $x$ in the canonical topology if and only if $x$ is \emph{strict}, meaning that every morphism $x' \to x$ is an isomorphism. Indeed, the pullback of the empty sieve along any morphism $x' \to x$ is the empty sieve on $x'$, which is a colimit sieve if and only if $x'$ is initial. For example, $\varnothing$ is a strict initial object of $\Set$ and of $\Open(X)$, but the zero group is not a strict initial object of $\Ab$, since every abelian group admits a morphism to it.\end{draft}
\end{example}

\begin{remark}\label{remark:chap:sheaves:sec:sites:colimit sieves need not be universal}
    A colimit sieve need not be a universal colimit sieve. For example, in $\Top$, the sieve generated by a single continuous map $q : Y \to X$ is a colimit sieve if and only if $q$ is a quotient map, and its pullback along a map $g : X' \to X$ is the sieve generated by the pullback $Y \times_X X' \to X'$. Quotient maps are not stable under pullback; see \Cref{exercise:chap:sheaves:sec:sites:colimit sieves in Top}.
\end{remark}

We saw above that a Grothendieck topology is subcanonical if and only if all of its covering sieves are universal colimit sieves. So if the collection of all universal colimit sieves formed a Grothendieck topology, it would automatically be the largest subcanonical topology. It is not obvious from the definitions that this is the case, but it is true! This largest Grothendieck topology is called the \emph{canonical topology} on $\catC$.

It's worth mentioning first that, just as for the collection of sieves on a given object, the collection of Grothendieck topologies on a given category forms a complete lattice under containment (any nonempty intersection of Grothendieck topologies is again a Grothendieck topology, and there is a largest Grothendieck topology on $\catC$ given by declaring every sieve to be covering). However, the join operation in this lattice is not easy to describe explicitly, so it is not always clear that there exists a largest Grothendieck topology satisfying some given property $P$. Instead, we can break up the problem of constructing the canonical topology into two steps: first, we produce the largest topology for which a given representable presheaf is a sheaf, then we form the intersection of all of these topologies.


\begin{proposition}\label{prop:chap:sheaves:sec:sites:largest topology for a presheaf}
    Let $\cal{F}$ be a presheaf on $\catC$. For each $x \in \catC$, let $\tau_{\cal{F}}(x)$ be the collection of sieves $S$ on $x$ such that $\cal{F}$ is local with respect to $g^* S$ for every morphism $g : x' \to x$. Then $\tau_{\cal{F}}$ is a Grothendieck topology on $\catC$. Moreover, $\cal{F}$ is a sheaf for $\tau_{\cal{F}}$, and $\tau_{\cal{F}}$ is the largest such topology: if $\tau$ is any Grothendieck topology for which $\cal{F}$ is a sheaf, then $\tau \subseteq \tau_{\cal{F}}$.
\end{proposition}
\begin{proof}
    First, assuming that $\tau_{\cal{F}}$ is indeed a Grothendieck topology, it is clear that $\cal{F}$ is a sheaf with respect to $\tau_{\cal{F}}$, since by definition $\cal{F}$ is local with respect to every covering sieve in $\tau_{\cal{F}}$. Moreover, if $\tau$ is any Grothendieck topology for which $\cal{F}$ is a sheaf, then for any $S \in \tau(x)$ and any $g : x' \to x$, we have $g^* S \in \tau(x')$ by (T1), and so $\cal{F}$ is local with respect to $g^* S$ by assumption. Hence $S \in \tau_{\cal{F}}(x)$, and so $\tau(x) \subseteq \tau_{\cal{F}}(x)$.

    Now we verify that $\tau_{\cal{F}}$ is indeed a Grothendieck topology. Axiom (T3) is straightforward, since $\cal{F}$ is local with respect to $g^* y(x) = y(x')$ for all $g : x' \to x$.

    Axiom (T1) is also immediate: for $S \in \tau_{\cal{F}}(x)$ and $f : x' \to x$ any morphism, we have $f^* S \in \tau_{\cal{F}}(x')$ because every pullback of $f^* S$ is also a pullback of $S$.

    For (T2), let $S \in \tau_{\cal{F}}(x)$, and let $T$ be a sieve on $x$ such that $f^* T \in \tau_{\cal{F}}(\dom f)$ for every $f \in S$. We must show that $\cal{F}$ is local with respect to $g^* T$ for every $g : x' \to x$. It suffices to show that $\cal{F}$ is local with respect to $T$, since for any $g : x' \to x$, the sieves $g^* S$ and $g^* T$ satisfy the same hypotheses as $S$ and $T$. By \Cref{prop:chap:sheaves:sec:sites:sheaf condition as limit}, this means we must show that, for every compatible family $(t_k)_{k \in T}$ of elements $t_k \in \cal{F}(\dom k)$, there is a unique $a \in \cal{F}(x)$ such that $\cal{F}(k)(a) = t_k$ for all $k \in T$.

    \emph{Uniqueness.} Suppose $a, b \in \cal{F}(x)$ satisfy $\cal{F}(k)(a) = \cal{F}(k)(b)$ for all $k \in T$, and let $f \in S$. For every $h \in f^* T$ we have $f \circ h \in T$, so
    \[\cal{F}(h)(\cal{F}(f)(a)) = \cal{F}(f \circ h)(a) = \cal{F}(f \circ h)(b) = \cal{F}(h)(\cal{F}(f)(b)).\]
    Since $\cal{F}$ is local with respect to $f^* T$, it follows that $\cal{F}(f)(a) = \cal{F}(f)(b)$. As this holds for every $f \in S$, and $\cal{F}$ is local with respect to $S$, we conclude that $a = b$.

    \emph{Existence.} Let $(t_k)_{k \in T}$ be a compatible family. For each $f \in S$, the family $(t_{f \circ h})_{h \in f^* T}$ is compatible, so since $\cal{F}$ is local with respect to $f^* T$, there is a unique $a_f \in \cal{F}(\dom f)$ with $\cal{F}(h)(a_f) = t_{f \circ h}$ for all $h \in f^* T$. The family $(a_f)_{f \in S}$ is compatible: if $f' = f \circ g$ with $f, f' \in S$, then $g \circ h \in f^* T$ for every $h \in f'^* T$, and so
    \[\cal{F}(h)(\cal{F}(g)(a_f)) = \cal{F}(g \circ h)(a_f) = t_{f \circ g \circ h} = t_{f' \circ h} = \cal{F}(h)(a_{f'})\]
    for all $h \in f'^* T$; since $\cal{F}$ is local with respect to $f'^* T$, this gives $\cal{F}(g)(a_f) = a_{f'}$. Since $\cal{F}$ is local with respect to $S$, there is some $a \in \cal{F}(x)$ with $\cal{F}(f)(a) = a_f$ for all $f \in S$. Finally, let $k \in T$ be arbitrary, and then let $h \in k^* S$ be arbitrary. We have $k \circ h \in T$ because $T$ is a sieve, so $\id_{\dom h} \in (k \circ h)^* T$. Thus,
    \[t_{k \circ h} = \cal{F}(\id_{\dom h})(a_{k \circ h}) = a_{k \circ h}.\]
    Then since $k \circ h \in S$, we have
    \[\cal{F}(h)(\cal{F}(k)(a)) = \cal{F}(k \circ h)(a) = a_{k \circ h} = t_{k \circ h} = \cal{F}(h)(t_k).\]
    Since $h \in k^* S$ was arbitrary, and $\cal{F}$ is local with respect to $k^* S$, we conclude that $\cal{F}(k)(a) = t_k$. Since $k \in T$ was arbitrary, we are done.
\end{proof}
\begin{theorem}\label{thm:chap:sheaves:sec:sites:canonical topology}
    Let $\catC$ be a category. The universal colimit sieves form a Grothendieck topology on $\catC$, called the \emph{canonical topology} (denoted $\tau_{\mathrm{can}}$). It is the largest subcanonical topology on $\catC$: a Grothendieck topology $\tau$ on $\catC$ is subcanonical if and only if $\tau \subseteq \tau_{\mathrm{can}}$.
\end{theorem}
\begin{proof}
    By \Cref{prop:chap:sheaves:sec:sites:when representables are sheaves}, a sieve $S$ on $x$ is a universal colimit sieve if and only if every representable presheaf is local with respect to $g^* S$ for every $g : x' \to x$; that is, if and only if $S \in \tau_{y(c)}(x)$ for every $c \in \catC$, in the notation of \Cref{prop:chap:sheaves:sec:sites:largest topology for a presheaf}. So the universal colimit sieves are precisely the sieves of $\bigcap_{c \in \catC} \tau_{y(c)}$, which is a Grothendieck topology by \Cref{prop:chap:sheaves:sec:sites:largest topology for a presheaf}. The final statement is \Cref{cor:chap:sheaves:sec:sites:subcanonical iff colimit sieves}, and applying it to the canonical topology itself shows that the canonical topology is subcanonical.
\end{proof}

\begin{proposition}\label{prop:chap:sheaves:sec:sites:examples are canonical}
    The Grothendieck topologies of \Cref{example:canonical topology on OpenX} and \Cref{example:canonical topology on Set} are the canonical topologies on $\Open(X)$ and on $\Set$, respectively. Moreover, in both $\Open(X)$ and $\Set$, every colimit sieve is universal.
\end{proposition}
\begin{proof}
    First consider $\Open(X)$. Since $\Open(X)$ is a poset, a cocone under a diagram in $\Open(X)$ is an upper bound for the objects in the diagram, and a colimit is a least upper bound. Least upper bounds in $\Open(X)$ are unions, so a sieve $S$ on $U$ is a colimit sieve if and only if $\bigcup S = U$, i.e.\ if and only if $S$ is generated by an open cover of $U$ (namely, by $S$ itself). Such sieves are stable under pullback: if $V \subseteq U$ is open, then $V^* S$ consists of the members of $S$ which are contained in $V$, and since $S$ is closed under passing to smaller open sets,
    \[\bigcup V^* S \supseteq \bigcup_{W \in S} (W \cap V) = V \cap \bigcup S = V.\]
    So every colimit sieve on $\Open(X)$ is universal, and the universal colimit sieves are exactly the sieves generated by open covers.

    Now consider $\Set$, and let $S$ be a sieve on a set $X$. Suppose first that $S$ is jointly surjective, and let $(s_f : \dom f \to Z)_{f \in S}$ be a cocone. Define $h : X \to Z$ by $h(x) := s_f(y)$ for any $f \in S$ and $y \in \dom f$ with $f(y) = x$. This is well-defined: if also $f'(y') = x$, then $i_x = f \circ i_y = f' \circ i_{y'}$ lies in $S$, so compatibility of the cocone gives $s_f(y) = s_{i_x}(\ast) = s_{f'}(y')$. By construction $h \circ f = s_f$ for all $f \in S$, and $h$ is unique with this property because $S$ is jointly surjective. Conversely, if $S$ is a colimit sieve, it is clearly jointly surjective. Finally, jointly surjective sieves are stable under pullback: if $g : X' \to X$ and $x' \in X'$, choose $f \in S$ and $y \in \dom f$ with $f(y) = g(x')$; then $g \circ i_{x'} = f \circ i_{y}$ lies in $S$, so $i_{x'} \in g^* S$, and its image contains $x'$. So every colimit sieve on $\Set$ is universal, and the universal colimit sieves are exactly the jointly surjective sieves.
\end{proof}

\begin{draft}
    In \Cref{prop:chap:sheaves:sec:sites:examples are canonical}, we saw that every colimit sieve in $\Open(X)$ or in $\Set$ is universal, in contrast with the situation in $\Top$ (\Cref{remark:chap:sheaves:sec:sites:colimit sieves need not be universal}). We now show that the same is true in the category of sheaves on any small subcanonical site, and we describe its canonical topology explicitly. We will also see that sheaves on such a category, for its canonical topology, are nothing new: they are the same thing as sheaves on the original site.

    For the rest of this subsection, fix a small site $(\catC,\tau)$ such that $\tau$ is subcanonical, and write $\cal{E} := \Sh_{\tau}(\catC)$. Since $\tau$ is subcanonical, every representable presheaf $y(c)$ is an object of $\cal{E}$, and the Yoneda embedding restricts to a fully faithful functor $y : \catC \to \cal{E}$. By the Yoneda lemma (\Cref{prop:Yoneda}), a section $a \in \cal{F}(c)$ of a sheaf $\cal{F} \in \cal{E}$ is the same thing as a morphism $a : y(c) \to \cal{F}$ in $\cal{E}$, and we will use the same letter for both. Under this identification, the restriction $\cal{F}(g)(a)$ of $a$ along a morphism $g : d \to c$ is the composite $a \circ y(g)$, and a morphism $f : \cal{F} \to \cal{G}$ in $\cal{E}$ sends the section $a$ to the section $f \circ a$.

    \begin{definition}\label{def:chap:sheaves:sec:sites:locally surjective sieve}
        Let $S$ be a sieve on an object $\cal{F} \in \cal{E}$. For each object $c \in \catC$ and each section $a \in \cal{F}(c)$, let $S_a$ be the sieve on $c$ in $\catC$ given by
        \[S_a := \{g : d \to c : a \circ y(g) \in S\}\]
        (this is a sieve because $S$ is). We say that $S$ is \emph{locally surjective} if $S_a \in \tau(c)$ for every object $c \in \catC$ and every section $a \in \cal{F}(c)$.
    \end{definition}

    \begin{remark}\label{remark:chap:sheaves:sec:sites:locally surjective families}
        Suppose that $S$ is generated by a family of morphisms $\{f_i : \cal{G}_i \to \cal{F}\}_i$. A morphism $y(d) \to \cal{F}$ lies in $S$ if and only if it factors through some $f_i$, so a morphism $g : d \to c$ lies in $S_a$ if and only if the section $\cal{F}(g)(a)$ lies in the image of $(f_i)_d : \cal{G}_i(d) \to \cal{F}(d)$ for some $i$. So $S$ is locally surjective if and only if, for every section $a \in \cal{F}(c)$, the morphisms $g : d \to c$ such that $\cal{F}(g)(a)$ lies in the image of one of the $f_i$ form a covering sieve on $c$. Informally: every section of $\cal{F}$ lies \emph{locally} in the image of one of the $f_i$. For sheaves on a topological space, this is a condition on stalks (\Cref{exercise:chap:sheaves:sec:sites:canonical topology on sheaves on a space}).
    \end{remark}

    \begin{theorem}\label{thm:chap:sheaves:sec:sites:canonical topology on a category of sheaves}
        Let $S$ be a sieve on $\cal{F} \in \cal{E}$. The following are equivalent:
        \begin{enumerate}[label=(\roman*)]
            \item $S$ is locally surjective;
            \item $S$ is a colimit sieve;
            \item $S$ is a universal colimit sieve.
        \end{enumerate}
        In particular, the covering sieves of the canonical topology $\tau_{\mathrm{can}}$ on $\cal{E}$ are exactly the locally surjective sieves, and every colimit sieve in $\cal{E}$ is universal.
    \end{theorem}
    \begin{proof}
        \emph{(i) implies (iii).} First, local surjectivity is stable under pullback: if $g : \cal{F}' \to \cal{F}$ is a morphism in $\cal{E}$ and $a \in \cal{F}'(c)$, then a morphism $h : d \to c$ lies in $(g^* S)_a$ if and only if $g \circ a \circ y(h) \in S$, so $(g^* S)_a = S_{g \circ a}$. So it suffices to show that every locally surjective sieve $S$ is a colimit sieve. Let $(s_f : \dom f \to \cal{H})_{f \in S}$ be a cocone in $\cal{E}$; we must show that there is a unique morphism $h : \cal{F} \to \cal{H}$ such that $h \circ f = s_f$ for all $f \in S$.

        Fix a section $a \in \cal{F}(c)$. For each $g : d \to c$ in $S_a$, we have $a \circ y(g) \in S$, so $s_{a \circ y(g)} : y(d) \to \cal{H}$ is a section of $\cal{H}$ over $d$. These sections form a compatible family: if $k : e \to d$, then $y(k)$ is a morphism from $a \circ y(g \circ k)$ to $a \circ y(g)$ in $S/\cal{F}$, so the cocone condition says that $s_{a \circ y(g)} \circ y(k) = s_{a \circ y(g \circ k)}$. Since $S_a \in \tau(c)$ and $\cal{H}$ is a sheaf, there is a unique section $h_a \in \cal{H}(c)$ such that
        \[h_a \circ y(g) = s_{a \circ y(g)} \quad \text{for all } g \in S_a.\]
        If $k : c' \to c$, then $S_{a \circ y(k)} = k^* S_a$, and $h_a \circ y(k)$ satisfies the defining property of $h_{a \circ y(k)}$, since $(h_a \circ y(k)) \circ y(g') = h_a \circ y(k \circ g') = s_{a \circ y(k) \circ y(g')}$ for all $g' \in k^* S_a$. Hence $h_{a \circ y(k)} = h_a \circ y(k)$, which says exactly that the maps $a \mapsto h_a$ assemble into a morphism $h : \cal{F} \to \cal{H}$ with $h \circ a = h_a$ for every section $a$.

        Now let $f : \cal{G} \to \cal{F}$ be a member of $S$, and let $b \in \cal{G}(c)$. Then $f \circ b \in S$, so $\id_c \in S_{f \circ b}$, and therefore
        \[h \circ f \circ b = h_{f \circ b} = s_{f \circ b} = s_f \circ b,\]
        where the last equality is the cocone condition for the morphism $b$ from $f \circ b$ to $f$ in $S/\cal{F}$. As $b$ was an arbitrary section of $\cal{G}$, we conclude that $h \circ f = s_f$. Finally, if $h' : \cal{F} \to \cal{H}$ also satisfies $h' \circ f = s_f$ for all $f \in S$, then for every section $a \in \cal{F}(c)$ and every $g \in S_a$ we have $(h' \circ a) \circ y(g) = s_{a \circ y(g)}$, so $h' \circ a = h_a = h \circ a$ by the uniqueness of $h_a$. Hence $h' = h$.

        \emph{(iii) implies (ii).} Every universal colimit sieve is a colimit sieve.

        \emph{(ii) implies (i).} Suppose that $S$ is a colimit sieve, and define a subpresheaf $\cal{I} \subseteq \cal{F}$ by
        \[\cal{I}(c) := \{a \in \cal{F}(c) : S_a \in \tau(c)\}.\]
        This is a subpresheaf by (T1), since $S_{a \circ y(k)} = k^* S_a$ for all $k : c' \to c$. It is moreover a sheaf: if $R \in \tau(c)$ and $(a_g)_{g \in R}$ is a compatible family of sections of $\cal{I}$, then since $\cal{F}$ is a sheaf, there is a unique $a \in \cal{F}(c)$ with $a \circ y(g) = a_g$ for all $g \in R$, and $a \in \cal{I}(c)$ by (T2), since $g^* S_a = S_{a \circ y(g)} = S_{a_g} \in \tau(\dom g)$ for all $g \in R$. Every member $f : \cal{G} \to \cal{F}$ of $S$ factors through $\cal{I}$: for every $b \in \cal{G}(c)$ we have $f \circ b \in S$, so $S_{f \circ b}$ is the maximal sieve, which is covering by (T3). Write $f = \iota \circ \bar{f}$, where $\iota : \cal{I} \hookrightarrow \cal{F}$ is the inclusion. Since $\iota$ is a monomorphism, the morphisms $\bar{f} : \dom f \to \cal{I}$ form a cocone, so since $S$ is a colimit sieve, there is a morphism $r : \cal{F} \to \cal{I}$ with $r \circ f = \bar{f}$ for all $f \in S$. Then $(\iota \circ r) \circ f = f = \id_{\cal{F}} \circ f$ for all $f \in S$, so the uniqueness part of the universal property of the colimit gives $\iota \circ r = \id_{\cal{F}}$. Hence $\cal{I} = \cal{F}$, i.e.\ $S$ is locally surjective.

        The final statement now follows from \Cref{thm:chap:sheaves:sec:sites:canonical topology}.
    \end{proof}

    For example, let $\catC = \ast$ be the terminal category, with the topology whose only covering sieve is the maximal one. Then $\cal{E} = \PSh(\ast) \simeq \Set$, and the locally surjective sieves are exactly the jointly surjective ones; so \Cref{thm:chap:sheaves:sec:sites:canonical topology on a category of sheaves} recovers the description of the canonical topology on $\Set$ from \Cref{prop:chap:sheaves:sec:sites:examples are canonical} (see \Cref{exercise:chap:sheaves:sec:sites:canonical topology on presheaves}).

    \begin{remark}\label{remark:chap:sheaves:sec:sites:locally surjective vs jointly epimorphic}
        The uniqueness argument in the proof of (i)~$\Rightarrow$~(iii) shows that the members of a locally surjective sieve $S$ on $\cal{F}$ are \emph{jointly epimorphic} in $\cal{E}$: if $h, h' : \cal{F} \to \cal{H}$ satisfy $h \circ f = h' \circ f$ for all $f \in S$, then $h = h'$. The converse also holds, but proving it requires a way to construct sheaves to test against, such as sheafification.
    \end{remark}

    Next, we show that sheaves on $\cal{E}$ for the topology $\tau_{\mathrm{can}}$ are determined by their values on the representable sheaves $y(e)$. The key point is that every sheaf is ``covered by representables'': for $\cal{K} \in \cal{E}$, let $T_{\cal{K}}$ be the sieve on $\cal{K}$ generated by all sections $z : y(e) \to \cal{K}$ of $\cal{K}$. For every section $a$ of $\cal{K}$, the sieve $(T_{\cal{K}})_a$ is maximal, so $T_{\cal{K}}$ is locally surjective, and hence it is a covering sieve for $\tau_{\mathrm{can}}$ by \Cref{thm:chap:sheaves:sec:sites:canonical topology on a category of sheaves}. For a sieve $S$ on $\cal{F} \in \cal{E}$, we also write $S|_{\catC} \subseteq \cal{F}$ for the subpresheaf consisting of the sections $a : y(c) \to \cal{F}$ which lie in $S$.

    \begin{proposition}\label{prop:chap:sheaves:sec:sites:canonical sheaves are determined on representables}
        Let $\cal{P} : \cal{E}^\op \to \Set$ be a sheaf for $\tau_{\mathrm{can}}$.
        \begin{enumerate}[label=(\alph*)]
            \item Let $\cal{K} \in \cal{E}$. If $p, p' \in \cal{P}(\cal{K})$ satisfy $\cal{P}(z)(p) = \cal{P}(z)(p')$ for every section $z : y(e) \to \cal{K}$ of $\cal{K}$, then $p = p'$.
            \item Let $S$ be a sieve on $\cal{F} \in \cal{E}$ such that every member of $S$ factors through some section $z : y(e) \to \cal{F}$ which lies in $S$. Then restricting compatible families $(p_k)_{k \in S}$ to the members of $S$ with representable domain gives a bijection
                  \[\cal{P}(S) \xrightarrow{\sim} \Hom_{\PSh(\catC)}(S|_{\catC}, \cal{P} \circ y).\]
            \item For every $\cal{K} \in \cal{E}$, the map
                  \[\cal{P}(\cal{K}) \to \Hom_{\PSh(\catC)}(\cal{K}, \cal{P} \circ y), \qquad p \mapsto \big(z \mapsto \cal{P}(z)(p)\big)\]
                  is a bijection, natural in $\cal{K}$ and in $\cal{P}$.
        \end{enumerate}
    \end{proposition}
    \begin{proof}
        (a) Every member of $T_{\cal{K}}$ has the form $z \circ v$ for some section $z : y(e) \to \cal{K}$, and then
        \[\cal{P}(z \circ v)(p) = \cal{P}(v)(\cal{P}(z)(p)) = \cal{P}(v)(\cal{P}(z)(p')) = \cal{P}(z \circ v)(p').\]
        So $p$ and $p'$ have the same image in $\cal{P}(T_{\cal{K}})$, and since $\cal{P}$ is local with respect to the covering sieve $T_{\cal{K}}$, we conclude that $p = p'$.

        (b) The map is injective: every member of $S$ can be written as $z \circ v$ with $z : y(e) \to \cal{F}$ in $S$, and then $p_{z \circ v} = \cal{P}(v)(p_z)$ for any compatible family $(p_k)_{k \in S}$. For surjectivity, let $\alpha : S|_{\catC} \to \cal{P} \circ y$ be a natural transformation. For each member $k : \cal{K} \to \cal{F}$ of $S$, choose a factorization $k = z \circ v$ with $z : y(e) \to \cal{F}$ in $S$, and set $p_k := \cal{P}(v)(\alpha(z)) \in \cal{P}(\cal{K})$. This does not depend on the choice of factorization: if $k = z_1 \circ v_1 = z_2 \circ v_2$, then for every section $w : y(e') \to \cal{K}$ we can write $v_i \circ w = y(m_i)$ for a unique $m_i : e' \to e_i$ (as $y$ is fully faithful), and $z_1 \circ y(m_1) = k \circ w = z_2 \circ y(m_2)$, so naturality of $\alpha$ gives
        \[\cal{P}(w)\big(\cal{P}(v_i)(\alpha(z_i))\big) = \cal{P}(y(m_i))(\alpha(z_i)) = \alpha(z_i \circ y(m_i)) = \alpha(k \circ w)\]
        for $i = 1, 2$; by (a), the two candidates for $p_k$ agree. The family $(p_k)_{k \in S}$ is compatible, since if $k = z \circ v$ then $k \circ u = z \circ (v \circ u)$, whence $p_{k \circ u} = \cal{P}(u)(p_k)$. Finally, taking the factorization $z = z \circ \id$ shows that $p_z = \alpha(z)$ for every $z \in S|_{\catC}$, so $(p_k)_{k \in S}$ maps to $\alpha$.

        (c) Since every section of $\cal{K}$ lies in $T_{\cal{K}}$, we have $T_{\cal{K}}|_{\catC} = \cal{K}$. The map in question is the composite of the bijection $\cal{P}(\cal{K}) \to \cal{P}(T_{\cal{K}})$ (recall that $T_{\cal{K}}$ is a covering sieve) with the bijection of (b) for $S = T_{\cal{K}}$. Naturality is straightforward to check.
    \end{proof}

    \begin{theorem}\label{thm:chap:sheaves:sec:sites:sheaves on a category of sheaves}
        The functors
        \begin{align*}
            \Sh(\cal{E},\tau_{\mathrm{can}}) & \to \Sh_{\tau}(\catC), & \cal{P} & \mapsto \cal{P} \circ y, \\
            \Sh_{\tau}(\catC) & \to \Sh(\cal{E},\tau_{\mathrm{can}}), & \cal{G} & \mapsto \Hom_{\cal{E}}({-},\cal{G}),
        \end{align*}
        are mutually inverse equivalences of categories.\footnote{Since $\cal{E}$ is not small, it is not clear a priori that the hom-sets of $\Sh(\cal{E},\tau_{\mathrm{can}})$ are sets. But they are: by \Cref{prop:chap:sheaves:sec:sites:canonical sheaves are determined on representables}(a), a morphism $\cal{P} \to \cal{Q}$ in $\Sh(\cal{E},\tau_{\mathrm{can}})$ is determined by its components at the objects $y(e)$ for $e \in \catC$.}
    \end{theorem}
    \begin{proof}
        For $\cal{G} \in \Sh_{\tau}(\catC)$, the presheaf $\Hom_{\cal{E}}({-},\cal{G})$ on $\cal{E}$ is a sheaf for $\tau_{\mathrm{can}}$, since the canonical topology is subcanonical (\Cref{thm:chap:sheaves:sec:sites:canonical topology}).

        Next, let $\cal{P} \in \Sh(\cal{E},\tau_{\mathrm{can}})$; we show that $\cal{P} \circ y$ is a sheaf on $(\catC,\tau)$. Let $R \in \tau(c)$, and let $\bar{R}$ be the sieve on $y(c)$ in $\cal{E}$ generated by the morphisms $y(g)$ for $g \in R$. Since $y$ is fully faithful and $R$ is a sieve, a morphism $y(h) : y(d) \to y(c)$ lies in $\bar{R}$ if and only if $h \in R$. In other words, $\bar{R}|_{\catC} = R$ as subpresheaves of $y(c)$, and $\bar{R}_h = h^* R$ for every section $h \in y(c)(d) = \Hom_{\catC}(d,c)$. By (T1), it follows that $\bar{R}$ is locally surjective, hence a covering sieve for $\tau_{\mathrm{can}}$ by \Cref{thm:chap:sheaves:sec:sites:canonical topology on a category of sheaves}. So $\cal{P}$ is local with respect to $\bar{R}$, and \Cref{prop:chap:sheaves:sec:sites:canonical sheaves are determined on representables}(b) gives bijections
        \[(\cal{P} \circ y)(c) = \cal{P}(y(c)) \xrightarrow{\sim} \cal{P}(\bar{R}) \xrightarrow{\sim} \Hom_{\PSh(\catC)}(R, \cal{P} \circ y) = (\cal{P} \circ y)(R)\]
        whose composite is the restriction map. That is, $\cal{P} \circ y$ is local with respect to $R$.

        Finally, the two composites are isomorphic to the identity functors. For $\cal{G} \in \Sh_{\tau}(\catC)$, we have $\Hom_{\cal{E}}(y({-}),\cal{G}) \cong \cal{G}$ by the Yoneda lemma. For $\cal{P} \in \Sh(\cal{E},\tau_{\mathrm{can}})$, \Cref{prop:chap:sheaves:sec:sites:canonical sheaves are determined on representables}(c) gives bijections
        \[\cal{P}(\cal{K}) \cong \Hom_{\PSh(\catC)}(\cal{K}, \cal{P} \circ y) = \Hom_{\cal{E}}(\cal{K}, \cal{P} \circ y),\]
        natural in $\cal{K} \in \cal{E}$ and in $\cal{P}$, where the equality holds because $\cal{E}$ is a full subcategory of $\PSh(\catC)$.
    \end{proof}

    When $\catC = \ast$, \Cref{thm:chap:sheaves:sec:sites:sheaves on a category of sheaves} says that every sheaf on $\Set$ for the canonical topology is representable. Together with \Cref{exercise:chap:sheaves:sec:sites:canonical topology on Set via subsets}, this recovers \Cref{thm:chap:sheaves:sec:intro:sheaves are representable}; in fact, the proof of \Cref{prop:chap:sheaves:sec:sites:canonical sheaves are determined on representables}(c) is a direct generalization of the proof of that theorem (see \Cref{exercise:chap:sheaves:sec:sites:canonical topology on presheaves}).

    \begin{remark}\label{remark:chap:sheaves:sec:sites:comparison lemma}
        We used the assumption that $\tau$ is subcanonical in order to regard the objects of $\catC$ as objects of $\cal{E}$, via the fully faithful functor $y$. The conclusion of \Cref{thm:chap:sheaves:sec:sites:sheaves on a category of sheaves} holds for every small site $(\catC,\tau)$ (compare the remark following \Cref{example:chap:sheaves:sec:spaces:non representable sheaf on Top}), but when $\tau$ is not subcanonical, the proof requires sheafification. \Cref{thm:chap:sheaves:sec:sites:sheaves on a category of sheaves} is also a special case of Verdier's \emph{comparison lemma} \cite[Exp.~III]{artin_theorie_1972} (see also \cite[\nopp C2.2]{johnstone_sketches_2002}), which says roughly the following: if $\catD$ is a site and $u : \catC \to \catD$ is a fully faithful functor from a small category $\catC$ such that every object of $\catD$ has a covering sieve generated by morphisms whose domains lie in the image of $u$, then restriction along $u$ gives an equivalence between sheaves on $\catD$ and sheaves on $\catC$ for the induced topology. In our situation, $\catD$ is $\cal{E}$ with the topology $\tau_{\mathrm{can}}$, and $u = y$; every $\cal{K} \in \cal{E}$ is covered by the sieve $T_{\cal{K}}$, and the induced topology on $\catC$ is $\tau$ (\Cref{exercise:chap:sheaves:sec:sites:recovering the topology from sheaves}).
    \end{remark}
\end{draft}

\subsection{Sheafification in Two Steps}

% TODO: Move the following draft material to after sheafification on sites has been constructed.
% It makes precise the first observation at the start of this subsection (sheaves as a non-free cocompletion).
% \begin{draft}
%     \begin{proposition}\label{prop:chap:sheaves:sec:sites:universal property of sheaves}
%         Let $(\catC,\tau)$ be a small site, and let $\cal{E}$ be a cocomplete category. Precomposition with the functor $\catC \to \Sh_{\tau}(\catC)$, $x \mapsto y(x)^{++}$, gives an equivalence between the category of cocontinuous functors $\Sh_{\tau}(\catC) \to \cal{E}$ and the full subcategory of $\Fun(\catC,\cal{E})$ spanned by the functors $F$ which \emph{send covering sieves to colimits}, meaning that for every covering sieve $S$ on every $x \in \catC$, the cocone $(F(f) : F(c) \to F(x))_{(f : c \to x) \in S}$ exhibits
%         \[F(x) \cong \colim_{(f : c \to x) \in S/x} F(c).\]
%     \end{proposition}
%     \begin{proof}[Sketch of proof]
%         By \Cref{prop:Yoneda}, a functor $F : \catC \to \cal{E}$ extends to a cocontinuous functor $\widehat{F} : \PSh(\catC) \to \cal{E}$, which has a right adjoint $e \mapsto \Hom_{\cal{E}}(F({-}),e)$. Since $S$ is the colimit of the representables $y(c)$ over $S/x$ (as observed after \Cref{prop:chap:sheaves:sec:sites:sheaf condition as limit}), the functor $\widehat{F}$ sends the inclusion $S \hookrightarrow y(x)$ to the comparison map $\colim_{S/x} F(c) \to F(x)$. It follows that $F$ sends covering sieves to colimits if and only if the right adjoint of $\widehat{F}$ takes values in sheaves, and this is exactly what is needed for $\widehat{F}$ to factor, uniquely up to isomorphism, through sheafification.
%     \end{proof}
%
%     \begin{proposition}\label{prop:chap:sheaves:sec:sites:subcanonical iff fully faithful}
%         Let $(\catC,\tau)$ be a small site. The functor $\catC \to \Sh_{\tau}(\catC)$, $x \mapsto y(x)^{++}$, is fully faithful if and only if $\tau$ is subcanonical.
%     \end{proposition}
%     \begin{proof}
%         For $x, c \in \catC$, the adjunction and the Yoneda lemma give
%         \[\Hom_{\Sh_{\tau}(\catC)}\left(y(x)^{++}, y(c)^{++}\right) \cong \Hom_{\PSh(\catC)}\left(y(x), y(c)^{++}\right) \cong y(c)^{++}(x),\]
%         and under this identification, the map $\Hom_\catC(x,c) \to \Hom\left(y(x)^{++},y(c)^{++}\right)$ induced by the functor is the component at $x$ of the unit $\eta_{y(c)} : y(c) \to y(c)^{++}$. So the functor is fully faithful if and only if $\eta_{y(c)}$ is an isomorphism for every $c$, i.e.\ if and only if every representable presheaf is a sheaf.
%     \end{proof}
%
%     This is a familiar phenomenon from presentations of algebraic objects: if we present an object by generators and relations, the generators map injectively into it exactly when the relations already held among them. Here the ``generators'' are the objects and morphisms of $\catC$, and the ``relations'' say that covering sieves become colimits. When $\tau$ is subcanonical, these relations are already true in $\catC$, so imposing them costs nothing, and $\catC$ embeds as a full subcategory of $\Sh_{\tau}(\catC)$. In this case \Cref{prop:chap:sheaves:sec:sites:universal property of sheaves} really does say that cocontinuous functors out of $\Sh_{\tau}(\catC)$ are the same as functors out of $\catC$ which \emph{preserve} the colimits designated by $\tau$, and we say that $\Sh_{\tau}(\catC)$ is a \emph{fully faithful} (not necessarily free) cocompletion of $\catC$, just as in \Cref{prop:chap:sheaves:sec:spaces:category of sheaves is cocomplete}. When $\tau$ is not subcanonical, the relations are genuinely new, and imposing them can collapse $\catC$.
%
%     In \Cref{example:chap:sheaves:sec:sites:dense topology}, the relation imposed by the sieve $S$ says that the induced map $y(\RR \setminus \{0\})^{++} \to y(\RR)^{++}$ is an isomorphism, even though $\RR \setminus \{0\} \hookrightarrow \RR$ is not an isomorphism in $\Open(\RR)$: imposing the relations has identified two distinct objects of $\Open(\RR)$.
%
%     \begin{remark}\label{remark:chap:sheaves:sec:sites:pullback stability and left exactness}
%         The condition (T1) plays a second role, on the side of the category of sheaves. If we imposed the relations coming from an arbitrary collection of sieves, we would still obtain a reflective subcategory of $\PSh(\catC)$ with a universal property like that of \Cref{prop:chap:sheaves:sec:sites:universal property of sheaves}. Stability of the relations under pullback is what guarantees that the resulting sheafification functor preserves finite limits, which will be essential later. So the word ``universal'' in \Cref{def:chap:sheaves:sec:sites:colimit sieve} and the axiom (T1) express the same requirement, seen from the two sides.
%     \end{remark}
% \end{draft}

%TODO \subsection{Points}

%TODO \subsection{Sheaves with Values in Abelian Categories}

%TODO \subsection{Closed Embeddings, Recollements, Artin Gluing}

\subsection*{Exercises}
\addcontentsline{toc}{subsection}{Exercises}

\begin{exercise}
    Show that the collection of sieves in \Cref{example:canonical topology on OpenX} satisfies the axioms of a Grothendieck topology.
\end{exercise}

\begin{exercise}
    Show that the collection of sieves in \Cref{example:canonical topology on Set} satisfies the axioms of a Grothendieck topology.
\end{exercise}

\begin{exercise}
    Check that \Cref{def:chap:sheaves:sec:sites:sheaf} recovers \Cref{def:chap:sheaves:sec:spaces:sheaf with values in algebraic category} when $(\catC, \tau) = (\Open(X), \tau)$ for a space $X$, where $\tau$ is the topology defined in \Cref{example:canonical topology on OpenX}.
\end{exercise}

\begin{exercise}\label{exercise:chap:sheaves:sec:sites:canonical topology on Set via subsets}
    Consider the Grothendieck topology on $\Set$ of \Cref{example:canonical topology on Set}.
    \begin{enumerate}[label=(\alph*)]
        \item Show that a sieve $S$ on a set $X$ is covering if and only if it is generated by the inclusions $X_i \hookrightarrow X$ of some family of subsets $\{X_i\}_i$ of $X$ with $\bigcup_i X_i = X$.
        \item Using \Cref{prop:chap:sheaves:sec:sites:sheaf condition as equalizer}, conclude that a presheaf $\cal{F}$ on $\Set$ is a sheaf for this topology if and only if it satisfies the sheaf condition $(\star)$ of \Cref{thm:chap:sheaves:sec:intro:sheaves are representable}. In particular, the sheaves on $\Set$ with respect to the canonical topology are exactly the representable presheaves.
    \end{enumerate}
\end{exercise}

\begin{exercise}
    For experts: if $\catD$ is an arbitrary category, we say that a presheaf $\cal{F} \in \PSh(\catC,\catD)$ is a sheaf if the $\Set$-valued presheaf
    \[\Hom_\catD(d,\cal{F}({-})) : \catC^\op \to \Set\]
    is a sheaf for all objects $d \in \catD$. Show that this recovers \Cref{def:chap:sheaves:sec:sites:sheaf} when $\catD$ is an algebraic category.
\end{exercise}

\begin{exercise}\label{exercise:chap:sheaves:sec:sites:dense topology is a topology}
    Show that the dense topology of \Cref{example:chap:sheaves:sec:sites:dense topology} satisfies the axioms of a Grothendieck topology.
\end{exercise}

\begin{exercise}\label{exercise:chap:sheaves:sec:sites:colimit sieves in Top}
    Let $q : Y \to X$ be a continuous map, and let $S$ be the sieve on $X$ in $\Top$ generated by $q$.
    \begin{enumerate}[label=(\alph*)]
        \item Show that $S$ is a colimit sieve if and only if $q$ is a quotient map.
        \item Show that for any continuous map $g : X' \to X$, the pullback sieve $g^* S$ is the sieve generated by the projection $Y \times_X X' \to X'$.
        \item Conclude that $S$ is a universal colimit sieve if and only if every pullback of $q$ is a quotient map. Find (or look up) a quotient map $q$ such that $q \times \id_{\mathbb{Q}}$ is not a quotient map, and conclude that not every colimit sieve in $\Top$ is universal.
    \end{enumerate}
\end{exercise}

\begin{draft}
    \begin{exercise}\label{exercise:chap:sheaves:sec:sites:canonical topology on sheaves on a space}
        Let $X$ be a topological space, and equip $\Open(X)$ with the topology $\tau$ of \Cref{example:canonical topology on OpenX}, which is subcanonical by \Cref{prop:chap:sheaves:sec:sites:examples are canonical} and \Cref{thm:chap:sheaves:sec:sites:canonical topology}. Thus $\Sh(X) = \Sh_{\tau}(\Open(X))$.
        \begin{enumerate}[label=(\alph*)]
            \item Show that a family of morphisms $\{f_i : \cal{G}_i \to \cal{F}\}_i$ in $\Sh(X)$ generates a locally surjective sieve if and only if, for every point $x \in X$, the maps of stalks $(f_i)_x : (\cal{G}_i)_x \to \cal{F}_x$ are jointly surjective.
            \item Using \Cref{thm:chap:sheaves:sec:spaces:sheaves are local homeos}, conclude that a sieve $S$ on an object $p : E \to X$ of $\Top^{\mathrm{lh}}/X$ is a covering sieve for the canonical topology on $\Top^{\mathrm{lh}}/X$ if and only if the images of the members of $S$ cover $E$.
        \end{enumerate}
    \end{exercise}

    \begin{exercise}\label{exercise:chap:sheaves:sec:sites:canonical topology on presheaves}
        Let $\catC$ be a small category, and let $\tau_{\mathrm{triv}}$ be the \emph{trivial topology} on $\catC$, whose only covering sieves are the maximal sieves.
        \begin{enumerate}[label=(\alph*)]
            \item Show that $\tau_{\mathrm{triv}}$ is a subcanonical Grothendieck topology, and that every presheaf on $\catC$ is a sheaf for $\tau_{\mathrm{triv}}$.
            \item Show that a sieve $S$ on $\cal{F} \in \PSh(\catC)$ is locally surjective (with respect to $\tau_{\mathrm{triv}}$) if and only if every section of $\cal{F}$ lies in the image of some member of $S$.
            \item Now let $\catC = \ast$ be the terminal category, so that $\PSh(\ast) \simeq \Set$. Show that (b) and \Cref{thm:chap:sheaves:sec:sites:canonical topology on a category of sheaves} recover the description of the canonical topology on $\Set$ in \Cref{prop:chap:sheaves:sec:sites:examples are canonical}. Then use \Cref{thm:chap:sheaves:sec:sites:sheaves on a category of sheaves} and \Cref{exercise:chap:sheaves:sec:sites:canonical topology on Set via subsets} to deduce \Cref{thm:chap:sheaves:sec:intro:sheaves are representable} again, and compare the proof of \Cref{prop:chap:sheaves:sec:sites:canonical sheaves are determined on representables}(c) with the proof of \Cref{thm:chap:sheaves:sec:intro:sheaves are representable}.
        \end{enumerate}
    \end{exercise}

    \begin{exercise}\label{exercise:chap:sheaves:sec:sites:recovering the topology from sheaves}
        Let $(\catC,\tau)$ be a small site with $\tau$ subcanonical, and let $R$ be a sieve on $c \in \catC$. Let $\bar{R}$ be the sieve on $y(c)$ in $\Sh_{\tau}(\catC)$ generated by the morphisms $y(g)$ for $g \in R$. Show that $R \in \tau(c)$ if and only if $\bar{R}$ is a covering sieve for the canonical topology on $\Sh_{\tau}(\catC)$. Conclude that $\tau$ is determined by the category $\Sh_{\tau}(\catC)$ together with the functor $y : \catC \to \Sh_{\tau}(\catC)$.\\
        \textit{Hint:} One direction is contained in the proof of \Cref{thm:chap:sheaves:sec:sites:sheaves on a category of sheaves}. For the other, consider $\bar{R}_{\id_c}$.
    \end{exercise}
\end{draft}